\scanpage{58}
% archive foliation: 57

\uline{Conclusions provisoires} (on pose $\overline{\mathbb{R}}=[-\infty,+\infty]$

\margin{\textcircled{1} On a $X',X''\in\mathcal{E}_{\chi}$ (resp. $\mathcal{E}_{\chi!}$) $\Longleftrightarrow X'\amalg X''\in\mathcal{E}_{\chi}$ (resp. $\mathcal{E}_{\chi!}$), et alors
\deleted{\unsure{résulte de c et e}}
\[
\chi(X)=\chi(X')+\chi(X'')\quad(\text{resp.}\ldots)
\]
résulte de c) et e) $(\varnothing\in\mathcal{E}_{\chi})$ resp. de b)}

\textcircled{\unsure{2}} \deleted{$\unsure{n\geq1}$}\inserted{$\forall n\in\mathbb{N}$}, $R^n,\overline{R}^n,B^n,S^n\in\mathcal{E}_{\chi}\cap\mathcal{E}_{\chi!}$, de plus
\[
\left\{\begin{aligned}
\chi(\unsure{R^n})&=\chi(\overline{R}^n)=\chi(B^n)=1\\
&\overset{\text{''}}{\chi_!}(\overline{R}^n)=\overset{\text{''}}{\chi_!}(B^n)\\
\chi_!(R^n)&=(-1)^n
\end{aligned}\right.
\qquad
\chi(S^n)=\chi_!(S^n)=1+(-1)^n
\]

\uline{Dém} $R^n,\overline{R}^n,B^n\sim pt$, donc par d) et e) ils sont $\in\mathcal{E}_{\chi}$ ; comme $\overline{R}^n,B^n$ compacts, par b) ils sont aussi $\in\mathcal{E}_{\chi!}$.

On a $S^n=S^n_+\cup S^n_-$ et $S^n_+\cap S^n_-\simeq S^{n-1}$, \inserted{Or $S^n_+\simeq S^n_-\simeq B^n$ donc $\in\mathcal{E}_{\chi}$} donc on conclut de c) par réc. sur $n$ que $S^n\in\mathcal{E}_{\chi}$ (\uline{NB} $S^0=\{pt\}\amalg\{pt\}\in\mathcal{E}_{\chi}$ par \unsure{d)|1)} et e), et $\chi(S^0)=1+1=2$), et on trouve
\[
\chi(S^n)=\chi(S^n_+)+\chi(S^n_-)-\chi(S^{n-1})\unsure{\ne|=}2-\chi(S^{n-1}),
\]
d'où par récurrence sur $n$ $\chi(S^n)=1+(-1)^n$. Comme $S^n$ cpt, $S^n\in\mathcal{E}_{\chi!}$, et $\chi_!(S^n)=\chi(S^n)$ par (a). Enfin, $R^n\simeq S^n-pt$, donc par (b) $R^n\in\mathcal{E}_{\chi!}$ et
\[
\chi(R^n)=\chi(S^n)-\chi(pt)=(1+(-1)^n)-1=(-1)^n,
\]
cqfd.

\textcircled{3} Soit $X_{-1}=\varnothing\unsure{\subset}X_0\subset X_1\subset\cdots\subset X_n=X$ filtration de l'espace $X$, tel que $\forall\,0\leq i\leq n$, $X_i-X_{i-1}$ \deleted{\unsure{est} $\unsure{\mathbb{R}}$} a un nb fini $f_i$ de comp. connexes ttes homéom. à $\mathbb{R}^i$ (i.e. $X_i-X_{i-1}\simeq f_i\times\mathbb{R}^i$). Alors $X\in\mathcal{E}_{\chi!}$, et on a
\[
\chi_!(X)=\sum(-1)^i f_i
\]
(donc si $X$ est cpt, on a $X\in\mathcal{E}_{\chi}$ et $\chi(X)=\sum(-1)^i f_i$ ...

\figure{\begin{tikzpicture}[x=1cm,y=1cm]
\draw (0,2)--(0,1.2)--(1.4,1)--(2.2,-.2)--(3.3,-.2);
\draw (0,1.5)--(3.4,1.5);
\end{tikzpicture}}
